% ===================================================================
%  calibrate-fr.tex — TOUTES les mesures du fac-simile « Livret explicatif des Tableaux Auxiliaires Delmas ».
%  Fichier PRODUIT par outils/kalibro.py : ne pas le modifier a la main.
%  Une seule constante y est physique — la largeur du papier — et elle
%  est passee en argument. Tout le reste vient de outils/inv-fr.json.
%
%  Releve : 65 pages de 35 lignes ou plus.
%    justification  1343 px  (ecart-type 182)
%    hauteur bloc   2679 px
%    pas des lignes 58.00 px  (ecart-type 7.67)
%    image de page  1686 x 2781 px
%    marge gauche   149 px
%    1re ligne a    65 px du bord
%  Echelle retenue : 15.450 px/mm, prise sur une hauteur de papier
%  de 180 mm (notice de bibliotheque) ; la page mesure alors
%  109.1 mm de large.
% ===================================================================

\newcommand{\VUlangue}{fr}
\newcommand{\VUmarque}{Gilles-Philippe Morin a concu, numerise et dirige la transcription de ce livre.}

% 1. GEOMETRIE
\newcommand{\VUpapierLargeur}{109.13mm}
\newcommand{\VUpapierHauteur}{180.00mm}
\newcommand{\VUtexteLargeur}{86.93mm}
\newcommand{\VUtexteHauteur}{173.40mm}
% LE FOLIO EST LA PREMIERE ENCRE DE LA PAGE, LE TEXTE VIENT APRES.
% Premiere version : les deux ordonnees etaient egales, et le folio se
% composait DANS la premiere ligne de texte — « 6 » barrait « esas » au
% folio 6. La mesure de outils/inventaire.py donne le haut du BLOC
% D'ENCRE, c'est-a-dire le folio sur toute page foliotee ; le corps du
% texte commence un rang plus bas. On ajoute donc un pas de ligne, et
% les pages sans folio (ouvertures de tableau) commencent a la meme
% ordonnee que les autres, comme au fac-simile.
\newcommand{\VUmargeSup}{7.96mm}
\newcommand{\VUfolioY}{4.21mm}

% 2. CORPS ET INTERLIGNAGE
\newcommand{\VUinterligne}{10.68pt}
\newcommand{\VUcorps}{10.70pt}
\newcommand{\VUcorpsNote}{8.45pt}
\newcommand{\VUinterligneNote}{9.03pt}
% Renfoncement d'alinea : provisoire, a relever (outils/mesures.py).
\newcommand{\VUalinea}{4.35mm}
\newcommand{\VUalineaNote}{3.48mm}
\newcommand{\VUblancAlineaValeur}{2.03pt}
\newcommand{\VUblancNote}{2.67pt}
\newcommand{\VUfiletnoteLargeur}{23.47mm}

% 3. ESPACE-MOT
% Largeur proche de celle de la fonte, elasticite large : chaque ligne
% trouve sa mesure sans que TeX ait a couper ailleurs qu'au point releve.
\newcommand{\VUespaceRatio}{0.32}
\newcommand{\VUespaceEtire}{0.30}
\newcommand{\VUespaceSerre}{0.14}

% 4. FONTES
\newcommand{\VUfonte}{XCharter-TLF}
\newcommand{\VUfonteGras}{ntxtlf}
\newcommand{\VUratioGras}{0.98}
